\chapter{Introduction}
\label{chap:introduction}

This chapter introduces the shipboard steelwork addressed by this project and explains why its design requires a traceable calculation process grounded in marine engineering practice. It outlines the findings from early exploratory testing, defines the engineering problem, and states the aim, objectives and research questions. It then explains the significance, scope and limitations of \NEMO{} and the organisation of this report.

\section{Background}
\label{sec:background}

Marine secondary steelwork connects equipment, services and appendages to the primary vessel structure. The term describes structural function rather than low consequence. A small attachment can carry the whole load of a suspended machine or restrain a pressurised pipe. Its failure can therefore interrupt operations or injure personnel even when the surrounding hull remains intact. This project considers lifting padeyes, pipe support brackets and stabilizer fin internal structures. These details share fabrication methods based on plates, welds and stock sections, but their loads and governing checks differ.

A lifting padeye transfers sling force through a shackle pin into a steel plate and its attachment welds. Yards use lifting points during machinery replacement, module installation and handling of fabricated assemblies. Some lifting points remain on board for maintenance; others support a particular construction operation. The pinhole interrupts the plate section and introduces local bearing pressure. The designer must consider tearing across the hole, fracture above it, shear-out, pin bearing, base bending and weld load transfer. Sling inclination changes the force even when the lifted mass stays constant. Unequal sharing between lifting points and sideways pull can change the design further. ASME BTH-1 supplies a named basis for pin-connected plate checks, while marine operations can impose requirements through DNV-ST-N001 or the applicable lifting code \cite{asmeBTH1,dnvN001,lrLifting2026}.

Pipe supports transfer the weight and inertia of piping into decks, bulkheads or other foundations. The supported mass includes the pipe metal, contents and insulation. The support must restrain the intended directions while allowing the movements required by the piping arrangement. A bracket can carry bending, shear and twisting because the pipe centre lies away from the supporting bulkhead. Its seat plate, knees, weld group and clamp form separate links in the load path. A strong bracket does not by itself establish a satisfactory support arrangement. The pipe can bend or vibrate between supports, and a flexible foundation can alter its response. MSS support practice and the structural checks in AISC 360 address different parts of this problem \cite{mssSP58,aisc2016}.

A roll-stabilizer fin generates a hydrodynamic force that produces a moment opposing vessel roll. The force passes through the fin skin, internal spars and ribs into a stock and its bearings. Ship speed, fin angle and planform affect the force. Stock position affects the twisting moment. The fin therefore connects fluid loading with structural sizing in a way that the other two parts do not. Perez and Blanke explain that roll-control performance also depends on wave excitation, control design and operating restrictions \cite{perez2012}. This project addresses structural sizing under a stated load envelope. It does not design a roll-control system or predict the reduction of roll for a vessel.

These details matter during both vessel operation and yard production. Plate selection affects cutting, handling and welding. A dimension that performs well in a continuous numerical search may require plate that the workshop does not hold. Thin details may also leave inadequate room for a shackle, clamp or practical weld. Conversely, copying a large attachment without checking the actual job can retain unnecessary steel and weld length. The useful engineering question is which available, buildable arrangement satisfies the stated checks, with enough information for another engineer to review the choice.

Engineers commonly address this question through hand calculations, spreadsheets and general computer-aided design and finite element analysis software. Hand calculations expose the load path and allow direct comparison with a rule. Spreadsheets repeat those calculations for different dimensions. CAD defines geometry and produces drawings or exchange models. FEA estimates stress and displacement where simple section models cannot describe the local field. Each method has value, but engineers must maintain agreement between their input values, units, geometry and assumptions. CAD and analysis integration research identifies this continuity as an engineering problem in its own right \cite{park2010,nolan2015}.

For example, changing a pinhole diameter affects net area, bearing clearance, shackle selection and the surface used to apply pressure in an analysis. Changing support spacing affects pipe weight on a bracket and the pipe natural frequency. Changing fin speed changes both the applied force and the pressure available before cavitation. Independent files can represent different versions of the same job unless the workflow preserves these relationships. An attractive stress contour cannot reveal all such disagreements. The report must show what was loaded and checked before it interprets the resulting stress.

\NEMO{} brings these calculations into a local desktop application built around shipboard engineering inputs. The engineer enters the service condition through structured questions and selects the plate thicknesses held by the shipyard. Part-specific rules screen candidate dimensions. The tool then builds exact solid geometry, runs the defined stress analysis and produces calculation working and fabrication files. Table~\ref{tab:intro-parts} summarises the three parts and the different marine engineering decisions that the shared workflow must support.

\begin{table}[H]
\centering
\caption{Parts covered by this project and their design basis}
\label{tab:intro-parts}
\small
\begin{tabularx}{\textwidth}{p{0.12\textwidth}p{0.16\textwidth}p{0.19\textwidth}XX}
\toprule
Part & Typical vessel location & Governing rule set used & Main loads & Outputs \\
\midrule
Lifting padeye & Machinery lifting point or deck-mounted assembly & ASME BTH-1 pin-connected plate and attachment checks & Sling force, in-plane force and sideways allowance & Plate, cheek and weld sizes; shackle selection; calculations; STEP and DXF \\
Pipe support bracket & Machinery-space or service-line bulkhead & AISC 360-16 ASD, with stated pipe-span practice & Pipe, fluid and insulation weight; motion inertia; friction & Seat, knee, weld and U-bolt sizes; span assessment; calculations; STEP and DXF \\
Stabilizer fin & Underwater hull appendage and stock region & IACS UR S10 provisions adapted as a structural basis & Quasi-static foil force, root bending and stock twisting & Skin, spar, rib and stock sizes; pressure and cavitation advice; calculations; STEP and rib DXF \\
\bottomrule
\end{tabularx}
\end{table}

\section{Findings from early exploratory testing}
\label{sec:early-testing}

Before fixing the architecture of \NEMO{}, exploratory prototype trials were conducted on a parameterised workflow for an aluminium equipment bracket, a steel lifting padeye and a stabilizer fin. Those preliminary tests combined first-order screening equations with a file-based bridge to Autodesk Fusion, which generated native geometry and exchange models. The trials showed that one software framework could handle structurally different marine details, while also exposing four practical engineering pitfalls that shaped the design requirements of \NEMO{}.

First, the analysis loop could not close automatically through the public Fusion API, which creates geometry and exchange files but does not expose Simulation study setup, solving or result retrieval \cite{autodeskStaticStress,autodeskAddins}. Engineers therefore had to transfer loads and boundary conditions into the finite element workspace by hand. Second, two manual Fusion bracket studies illustrated how fragile that handoff is. One study applied 1.472~N instead of the intended 1,472~N, and the other lacked a traceable applied-load record. Both runs completed without solver warnings, yet a thousand-fold load error produced artificially low stresses and high safety factors. Both studies had to be rejected, confirming that a completed stress plot cannot substitute for automatic load-transfer and equilibrium checks.

Third, the early screening formulas relied on unsourced tuning factors and assumed hydrodynamic coefficients rather than named marine design codes and derived foil loads. An optimiser can exploit an unverified screening expression to reduce mass while violating pin-connection, weld-throat or slender-plate limits. Finally, requiring Python environment setup and manual CAD configuration created an unnecessary barrier for routine marine design review. These early observations shifted the emphasis of the project from unconstrained mass minimisation toward code-based sizing over shipyard plate stock, automated open-source numerical verification and traceable calculation working.

\section{Problem statement}
\label{sec:problem-statement}

The engineering problem addressed by this project is the gap between selecting dimensions for shipboard secondary steelwork and producing a reproducible, code-traceable structural justification for those dimensions. General spreadsheets, CAD packages and finite element solvers supply the individual operations separately. In daily marine and shipyard practice, the engineer must still translate service conditions into load components, apply the relevant design standard, match dimensions to available plate stock and catalogue hardware, rebuild geometry, configure boundary conditions and compile a reviewable calculation sheet. Each manual transfer risks unit errors, mismatched geometry or omitted load cases, as seen during the early exploratory bracket trials.

A dependable marine design workflow must preserve the engineering basis from initial load derivation through final stress confirmation. It must distinguish published code provisions from stated yard practice, and required structural checks from advisory guidance. It must reject any candidate that violates a governing constraint regardless of low steel mass, and it must define the scope of its finite element model clearly. A converged stabilizer stock analysis does not qualify the surrounding hull bearings, just as bracket equilibrium does not verify the supporting bulkhead plating. Verification evidence must match the physical quantity on which the design decision rests \cite{asmeVV10,nasa7009}.

This project investigates how code-based structural rules, shipyard plate-stock sizing, automated three-dimensional finite element verification and traceable calculation-package generation can be unified into one desktop tool for lifting padeyes, machinery-space pipe brackets and roll-stabilizer fins. The engineering contribution lies in integrating and verifying established naval architecture and structural mechanics methods for these shipboard details, and in testing what actually governs their dimensions across realistic vessel operating envelopes.

\section{Aim and objectives}
\label{sec:aim-objectives}

The aim of this project is to develop and verify a rule-based marine structural design tool (\NEMO{}) that sizes lifting padeyes, machinery-space pipe support brackets and roll-stabilizer fin internal structures from available shipyard plate stock, confirms each design through automated three-dimensional finite element analysis, and issues a traceable engineering calculation package and fabrication files for review.

The objectives are:
\begin{enumerate}
\item Formulate and implement code-based structural check and stock-aware sizing models for three shipboard details, namely lifting padeyes to ASME BTH-1, pipe support brackets to AISC 360-16 with machinery-space pipe-span vibration screening, and roll-stabilizer fin internal structures to IACS UR S10 with independent foil-load and cavitation models, searching available shipyard plate thicknesses to deliver buildable designs with full hand-calculation working, STEP and DXF outputs, and a reviewable PDF calculation package through a local desktop interface.
\item Establish an automated open-source geometry and quadratic tetrahedral finite element pipeline (build123d on OpenCASCADE, Gmsh, scikit-fem and PyAMG) with three-level mesh refinement and automatic force and moment equilibrium checks, and verify the structural, geometric and hydrodynamic solvers against hand calculations, closed-form theory, published reference solutions and automated regression tests.
\item Evaluate the verified tool across representative vessel case studies and operating-envelope parameter sweeps (load, sling inclination and obliquity, pipe nominal size and support spacing, ship speed and stock bearing arrangement) to determine the structural, catalogue, vibration and hydrodynamic limits that govern each marine detail in practice.
\end{enumerate}

\section{Research questions}
\label{sec:research-questions}

\begin{enumerate}
\item How can rule-based structural checks, shipyard plate-stock constraints and traceable hand-calculation working be unified across structurally distinct marine details without obscuring code provisions or engineering assumptions?
\item How closely do the automated solid-geometry, quadratic finite element and potential-flow foil solvers agree with hand calculations, closed-form elasticity and beam theory, and published reference data?
\item Across realistic shipboard operating envelopes, which structural, catalogue, vibration and hydrodynamic limits govern lifting padeyes, machinery-space pipe brackets and stabilizer fins in practice?
\end{enumerate}

\section{Significance of the project}
\label{sec:significance}

The marine engineering value of this work lies in keeping the physical basis and code provenance attached to every selected dimension. Each check carries its governing equation, substituted values, numerical result and named code or yard practice. A reviewing engineer or surveyor can follow the load path from lifted mass, ship-motion acceleration or vessel speed straight to the governing weld, pin ligament, pipe resonance band or stock section without reading source code.

Stock-aware sizing connects numerical search with shipyard fabrication reality. Plate thicknesses are chosen from the yard's inventory rather than continuous optima that cannot be ordered or cut, and dimensions are rounded and rechecked alongside catalogue shackle jaws, U-bolt clearances and minimum fillet weld legs. Offering three reserve targets exposes the trade-off between steel mass, catalogue hardware steps and structural margin.

The automated open-source analysis pipeline removes the fragile manual transfer between calculation and finite element simulation. Built on OpenCASCADE, Gmsh and scikit-fem \cite{gmsh2009,skfem2020}, the solver enforces three mesh refinements and checks six-degree-of-freedom static equilibrium automatically on every design. A built-in verification suite lets any engineer rerun the reference benchmarks on the local machine before relying on a calculation package, keeping numerical verification distinct from physical proof testing.

Finally, packaging the workflow as a local browser application with live three-dimensional previews and complete calculation records removes programming prerequisites from routine marine sizing while retaining every engineering decision in the project folder. Formal user trials with practising engineers remain outside the present scope and are recorded as future work.

\section{Scope and limitations}
\label{sec:scope-limitations}

The scope covers three shipboard steelwork families. The lifting padeye consists of a rounded head, pinhole, base fillet welds and optional cheek plates. The machinery-space pipe support bracket consists of a seat plate, two triangular knees and a U-bolt, paired with a pipe-span natural frequency check against engine and propeller excitations. The roll-stabilizer fin retains a specified symmetric NACA foil profile while sizing skin plating, front and rear spars, internal ribs and a solid forged stock. The tool derives service loads, evaluates rule checks, builds STEP and DXF outputs, runs built-in FEA and compiles a PDF calculation package.

The finite element models use linear-elastic isotropic material behaviour and quadratic tetrahedral displacement elements. The padeye and bracket solids are analysed under idealised rigid-base restraints and surface tractions; the stabilizer FEA models the solid stock alone and checks it against combined bending and torsion beam theory, while the thin-walled fin box is checked by closed-section and plate formulas. Three mesh refinements, six-degree-of-freedom equilibrium and benchmark comparisons supply numerical verification evidence, not physical validation of service behaviour \cite{asmeVV10}.

Surrounding vessel structures lie outside the calculation boundary. \NEMO{} does not check deck plating or under-deck stiffeners beneath a padeye, bulkhead plating behind a pipe bracket, or hull trunks, bearings and stock-to-fin couplings around a stabilizer. A fixed support face models an assumed rigid boundary rather than certifying the ship structure behind it.

Fatigue life, nonlinear pin contact, fabrication imperfections and corrosion wastage are excluded. The pipe model omits thermal expansion, anchor loads and water hammer, and ship-motion accelerations are engineering inputs drawn from the vessel's applicable class rules rather than computed from hull geometry. Fin loads are quasi-static and omit dynamic fin actuation, wave slamming and unsteady hull interaction; the Hess-Smith panel method predicts inviscid pressure and cavitation onset rather than separated or cavitating flow. IACS UR S10 rudder rules are adapted as the nearest published structural basis for the fin and do not constitute class approval \cite{iacsS10}. The optional Autodesk Fusion connector builds parametric padeyes and provides a static stress cross-check sheet, but live physical testing and formal usability trials with external users remain outside this project.

The benchmarks and case studies support engineering claims only within these stated boundaries. \NEMO{} is a design aid; a competent marine engineer must review the load basis, applicable class rules, omitted failure modes and calculation package before fabrication or installation.

\section{Organisation of the report}
\label{sec:report-organisation}

Chapter~\ref{chap:literature} reviews the design rules, hydrodynamic models, numerical methods and software practices on which the project draws. Chapter~\ref{chap:methodology} describes the engineering models, sizing procedure, implementation and verification methods. Chapter~\ref{chap:results} presents the benchmark evidence, case studies and input sensitivity. Chapter~\ref{chap:workplan} sets out the project schedule, and Chapter~\ref{chap:budget} records the required resources and costs. Chapter~\ref{chap:conclusion} evaluates the objectives, states the supported conclusions and recommends further work.
